\section*{Course context}
\phantomsection
\makeatletter
\def\@currentlabel{Course context}\def\@currentlabelname{Course context}%
\def\cref@currentlabel{[context][1][]Course context}%
\makeatother
\crefformat{context}{#2#1#3}\Crefformat{context}{#2#1#3}
\label{sec:logistics}
\addcontentsline{toc}{section}{Course context}

This section places Lecture~2 within the course. It gives the question that motivates the course,
the four units of the semester, and what the exams assess. The question also explains why we treat
the model as one component of a larger system.

\subsection*{The question that motivates the course}

The course is built around the following question:
\begin{quote}
How do we place a \emph{probabilistic} model \emph{inside a software system} without surrendering
the \emph{guarantees} expected of production software?
\end{quote}
In this question, the \textbf{model} is a large language model, the component that generates a
reply to each request it receives. The table below explains why each phrase in italics belongs in
the question.

\begin{center}
\begin{adjustbox}{width=\linewidth}\begin{tabularx}{\linewidth}{@{}L{0.3\linewidth}X@{}}
  \toprule
  \textbf{Phrase} & \textbf{Why it is in the question} \\
  \midrule
  \textbf{Probabilistic} & The model samples each reply from a probability distribution, so two
    calls with the same input can return different outputs. Components that use the model's output
    must not assume that the outputs agree. \\
  \addlinespace
  \textbf{Inside a software system} & The model is one component among several. Other components
    call it, validate its output, and can refuse that output. \\
  \addlinespace
  \textbf{Guarantees} & A guarantee is a promise that a system makes to its callers: a behavior
    that a caller is entitled to expect and to rely on every time. \\
  \bottomrule
\end{tabularx}\end{adjustbox}
\end{center}

Together, the three phrases describe the problem that the course studies. A production system must
keep its guarantees on every call. The model inside it, by contrast, can answer the same request
differently each time. The course asks how the other components of the system, which call the model
and validate its output, can keep those guarantees.

The agent that we build around the model (\cref{sec:agents} defines the term) is also one component
of a larger system, which we call a \textbf{production agentic system}. By the end of the semester,
we will design such a system.

\subsection*{How the semester is organized}

The semester is divided into four units, each devoted to one part of a production agentic system.
The table below lists the units in order, with the weeks that each one takes, its subject, and the
part of the system that it studies.

\par\vskip 0pt plus 1fil\penalty500\vskip 0pt plus -1fil\relax
\begin{center}
\begin{adjustbox}{width=\linewidth}\begin{tabularx}{\linewidth}{@{}L{0.17\linewidth}L{0.34\linewidth}X@{}}
  \toprule
  \textbf{Unit (weeks)} & \textbf{Subject} & \textbf{What it studies} \\
  \midrule
  \textbf{1 (1--3)} & Agent architectures and control & The orchestrator: who chooses the next
    operation, and who decides to stop \\
  \addlinespace
  \textbf{2 (4--6)} & Tool interfaces and concurrency & The tools and the subagents: their
    contracts, their dependencies, and bounded fan-out (independent operations run in parallel,
    with a limit on how many run at the same time) \\
  \addlinespace
  \textbf{3 (8--9)} & Context, state, and persistence & Memory and context: what survives
    truncation of the context (the input that the model receives in one call), and what survives a
    restart \\
  \addlinespace
  \textbf{4 (10--12)} & Deployment, recovery, and observability & The service around the agent,
    its tools, and its memory: partial failure, retries, and traces \\
  \bottomrule
\end{tabularx}\end{adjustbox}
\end{center}

Lecture~2 belongs to Unit~1, agent architectures and control. Unit~1 studies the
\textbf{orchestrator}, the component that decides which operation runs next and when the agent
stops.

\subsection*{What the exams assess}

The exams are written and closed book. They assess conceptual understanding, architectural
reasoning, and the ability to analyze system behavior.
\bi
\item \textbf{Midterm.} The midterm, on October~23, lasts 40 minutes and covers Units~1 and~2.
\item \textbf{Final.} The final, on December~17, lasts 40 minutes and covers Units~1 to~4.
\ei
Because Lecture~2 belongs to Unit~1, both exams cover the material in these notes.
